% !TeX TXS-program:compile = txs:///lualatex | txs:///view
\documentclass{article}
\usepackage[T1]{fontenc}
\usepackage{libertine}
\usepackage[shortlabels]{enumitem}
\usepackage[
colorlinks=true,
urlcolor=blue,
linkcolor=black,
anchorcolor=black,
citecolor=black
]{hyperref}
\hypersetup{
	pdfinfo={
		Title={Clase 1 -- ejemplo 1},
		Author={Oromion},
		Keywords={TeX, TeXlive, lenguaje de marcado, Lamport, Knuth},
		Subject={Curso Vacacional Básico},
		Producer={TeXstudio 4.9.1},
		Creator={pdfTeX 3.141592653-2.6-1.40.27 (TeX Live 2026/dev/Arch Linux)}
	}
}
\begin{document}

Veamos la siguiente numeración de algunas ramas de la matemática.
\begin{enumerate}[label={\alph*)}]
	\item Análisis complejo.

	\item Teoría de colas.

	\item Geometría hiperbólica.

	\item Álgebra no conmutativa.
\end{enumerate}
\end{document}
